\paragrafos{Θετικοί και αρνητικοί αριθμοί: ο κόσμος κάτω από το μηδέν}{Α.7.1, «Θετικοί και αρνητικοί αριθμοί (ρητοί αριθμοί) – Η ευθεία των ρητών – Τετμημένη σημείου»}

\begin{anagke}
Η Α3 κάνει εκδρομή στο Μέτσοβο. Στο πούλμαν το ψηφιακό θερμόμετρο γράφει $-4$. «Τι είναι αυτό το μείον;» ρωτάει η Πολυξένη. «Τέσσερις βαθμούς κάτω από το μηδέν», απαντάει ο Νίκος. Ο Στέλιος ταράζεται: «Κάτω από το μηδέν; Το μηδέν δεν είναι το τίποτα; Πώς μπορεί να υπάρχει κάτι μικρότερο από το τίποτα;».

Στο ξενοδοχείο ο ανελκυστήρας έχει κουμπιά $-2$, $-1$, $0$, $1$, $2$, $3$. Ο Ηλίας πατάει το $-1$ και κατεβαίνει στο υπόγειο με το πάρκινγκ. Η Αγγελική βλέπει στον πίνακα του καιρού: Αθήνα $+7$, Ιωάννινα $-3$, Φλώρινα $-10$, Χανιά $+11$. Η Μυρτώ θυμάται ότι στο ταμείο της παρέας κάποτε έλειπαν $8$ ευρώ, ένα χρέος. Και στο ντοκιμαντέρ που είδαν, ο καρχαρίας βρισκόταν $200$ μέτρα κάτω από την επιφάνεια ενώ το αεροπλάνο πετούσε $200$ μέτρα πάνω από αυτήν.

Η κυρία Αλεξάνδρα το συνοψίζει: «Οι φυσικοί, οι δεκαδικοί και τα κλάσματα που ξέρετε μετράνε πόσα. Όταν όμως χρειάζεται να πούμε «πάνω ή κάτω», «κέρδος ή ζημιά», «πάνω ή κάτω από το μηδέν», χρειαζόμαστε και αριθμούς με κατεύθυνση». Κάθε «δεξιά» προϋποθέτει ένα «αριστερά» και κάθε «πάνω» ένα «κάτω»: ο κόσμος των αριθμών έχει δύο πλευρές.
\end{anagke}

\begin{lexi}{θετικός, αρνητικός, ρητός}{θετικός, αρνητικός, ρητός}
\lexiword{Θετικός} είναι αυτός που «τίθεται», που βάζουμε ως βάση. \lexiword{Αρνητικός} είναι αυτός που «αρνείται», που πηγαίνει ανάποδα από τον θετικό. Το \lexiword{πρόσημο} είναι το «προ-σήμα», το σημάδι που μπαίνει πριν από τον αριθμό. Ο \lexiword{ακέραιος} είναι «ακέραιος», ολόκληρος. Ο \lexiword{ρητός} (από το «ρητός», αυτός που μπορεί να ειπωθεί) είναι ο αριθμός που εκφράζεται ως κλάσμα, δηλαδή ως λόγος δύο ακεραίων.
\end{lexi}

\begin{metafrasi}
Τέσσερις βαθμοί κάτω από το μηδέν. & $-4$\\
Αεροπλάνο $200$ μ. πάνω από τη θάλασσα. & $+200$\\
Καρχαρίας $200$ μ. κάτω από την επιφάνεια. & $-200$\\
Χρέος $8$ ευρώ. & $-8$\\
Κέρδος $15$ ευρώ. & $+15$\\
Δεύτερο υπόγειο. & $-2$\\
\end{metafrasi}

\tmima{Η θεωρία}

Τα σύμβολα «$+$» και «$-$» λέγονται \textbf{πρόσημα}. Γράφονται πριν από τους αριθμούς και τους χαρακτηρίζουν, αντίστοιχα, ως \textbf{θετικούς} ή \textbf{αρνητικούς}.
$$+3,\ +\tfrac{3}{4},\ +15{,}7,\ +0{,}352\ \text{ είναι θετικοί αριθμοί},\qquad -2,\ -10,\ -28{,}95,\ -0{,}098\ \text{ είναι αρνητικοί}.$$
Όταν μιλάμε μόνο για θετικούς, μπορούμε να παραλείπουμε το $+$: γράφουμε $7$ αντί για $+7$. Το $0$ \textbf{δεν είναι ούτε θετικός ούτε αρνητικός}.

\begin{orismos}{ομόσημοι, ετερόσημοι, ακέραιοι, ρητοί}
Δύο αριθμοί είναι \textbf{ομόσημοι} όταν έχουν το ίδιο πρόσημο και \textbf{ετερόσημοι} όταν έχουν διαφορετικό. \textbf{Ακέραιοι} είναι οι φυσικοί μαζί με τους αντίστοιχους αρνητικούς: $\dots,\,-3,\,-2,\,-1,\,0,\,1,\,2,\,3,\dots$ \textbf{Ρητοί} είναι όλοι οι αριθμοί που ξέρουμε (φυσικοί, κλάσματα, δεκαδικοί) μαζί με τους αντίστοιχους αρνητικούς.
\end{orismos}

Για παράδειγμα οι $-7$, $-0{,}58$, $-\frac{3}{4}$ είναι ομόσημοι (όλοι αρνητικοί), οι $+1{,}25$, $+10$, $5$ επίσης, ενώ οι $-7$ και $+0{,}58$ είναι ετερόσημοι.

\tmimasmall{Μια ιστορική σημείωση}

Οι αρνητικοί αριθμοί εμφανίζονται για πρώτη φορά σε ένα κινέζικο μαθηματικό βιβλίο, τα «Μαθηματικά σε εννέα βιβλία», που τοποθετείται στην εποχή της δυναστείας των Χαν (206 π.Χ. -- 220 μ.Χ.). Στην ευρωπαϊκή παράδοση ο Διόφαντος (γύρω στο 250 μ.Χ.) τους χρησιμοποιεί στους ενδιάμεσους υπολογισμούς του, αλλά ως λύση ενός προβλήματος δέχεται πάντα θετικό αριθμό. Για πολλούς αιώνες οι μαθηματικοί τους θεωρούσαν «παράλογους» ή «ψεύτικους» αριθμούς. Σήμερα τους χρησιμοποιούμε στο θερμόμετρο, στον λογαριασμό της τράπεζας, στο υψόμετρο.

\tmimasmall{Η ευθεία των ρητών και η τετμημένη}

Στον ημιάξονα $Ox$ των αριθμών τοποθετούσαμε τους φυσικούς, δεξιά από την αρχή $O$. Αν προεκτείνουμε προς τα αριστερά (στον αντικείμενο ημιάξονα $Ox'$) με την ίδια μονάδα, παίρνουμε τον \textbf{άξονα $x'x$}, την ευθεία των ρητών. Τα σημεία δεξιά του $O$ παριστάνουν θετικούς και τα αριστερά αρνητικούς. Ο αριθμός που αντιστοιχεί σε ένα σημείο λέγεται \textbf{τετμημένη} του σημείου και γράφουμε $A(-3)$.

\begin{center}
\akgr{-5}{5}{-3/A,2.5/B,-1.5/\Gamma,4/\Delta}
\end{center}

Στο σχήμα: $A(-3)$, $B(2{,}5)$, $\Gamma(-1{,}5)$, $\Delta(4)$. Στην ευθεία των ρητών οι αρνητικοί βρίσκονται αριστερά του $0$ και, όσο πηγαίνουμε αριστερά, οι αριθμοί μικραίνουν.

\begin{prosoxi}
\begin{itemize}[leftmargin=1.2em,itemsep=0.3ex,topsep=2pt]
\item Το $-3$ δεν είναι «τρία ευρώ λιγότερο από το τίποτα». Είναι μια θέση, τρεις μονάδες αριστερά του $0$, και η μονάδα μπορεί να είναι βαθμός, ευρώ ή μέτρο.
\item Το $0$ δεν έχει πρόσημο. Δεν γράφουμε $+0$ ή $-0$.
\item Στη μονάδα του άξονα πρέπει οι αποστάσεις να είναι ίσες: το $-1{,}5$ είναι στη μέση του $-2$ και του $-1$.
\item Μην μπερδεύεις το πρόσημο με την πράξη: στο $5-3$ το μείον είναι αφαίρεση, στο $-3$ είναι πρόσημο.
\end{itemize}
\end{prosoxi}

\tmima{Λυμένα παραδείγματα}

\paradeigma{Αριθμοί για μεγέθη}
\tixero Οι πόλεις και οι θερμοκρασίες: Αθήνα $7$° πάνω από το $0$, Ιωάννινα $3$° κάτω, Φλώρινα $10$° κάτω. Επίσης ένα ντεπόζιτο είναι $12$ ευρώ χρέος.\par
\tithelo Να τα γράψω ως ρητούς.\par
\possindeo Πάνω από το μηδέν είναι θετικό και κάτω αρνητικό. Αθήνα $+7$, Ιωάννινα $-3$, Φλώρινα $-10$. Το χρέος είναι $-12$.\par
\apantisi $+7$, $-3$, $-10$ και $-12$.\par
\elegxos Η Φλώρινα ($-10$) είναι πιο κρύα από τα Ιωάννινα ($-3$): θα το συγκρίνουμε στην επόμενη παράγραφο.

\paradeigma{Ομόσημοι ή ετερόσημοι;}
\tixero Τα ζεύγη $(-3,\,-8)$, $(+4,\,-4)$, $(0{,}5,\,\frac{7}{2})$, $(-2,\,2{,}5)$.\par
\tithelo Να πω τι είναι.\par
\possindeo Κοιτάζω τα πρόσημα. $-3$ και $-8$ είναι αρνητικοί: ομόσημοι. $+4$ και $-4$: ετερόσημοι. $0{,}5$ και $\frac{7}{2}$ θετικοί: ομόσημοι. $-2$ και $2{,}5$: ετερόσημοι.\par
\apantisi Ομόσημοι: το πρώτο και το τρίτο. Ετερόσημοι: το δεύτερο και το τέταρτο.\par
\elegxos Δύο αριθμοί με το ίδιο πρόσημο βρίσκονται στην ίδια πλευρά του $0$ στον άξονα.

\paradeigma{Ακέραιοι και ρητοί}
\tixero Οι αριθμοί $-4$, $2{,}5$, $0$, $\frac{9}{3}$, $-\frac{1}{2}$, $10$.\par
\tithelo Ποιοι είναι ακέραιοι.\par
\possindeo Ακέραιοι είναι οι αριθμοί που δεν έχουν δεκαδικό ή κλασματικό μέρος. $\frac{9}{3}=3$ είναι ακέραιος. Όλοι είναι ρητοί.\par
\apantisi Ακέραιοι: $-4$, $0$, $\frac{9}{3}$, $10$. Οι $2{,}5$ και $-\frac{1}{2}$ είναι ρητοί αλλά όχι ακέραιοι.\par
\elegxos Κάθε ακέραιος είναι ρητός: π.χ. $-4=-\frac{4}{1}$.

\paradeigma{Σημεία στον άξονα}
\tixero Στο σχήμα της θεωρίας.\par
\tithelo Να βρω τετμημένες και να τοποθετήσω τον $-\frac{1}{2}$.\par
\possindeo Το $A$ είναι $3$ μονάδες αριστερά του $0$: $-3$. Το $B$ είναι στη μέση του $2$ και του $3$: $2{,}5$. Το $\Gamma$ στη μέση του $-2$ και $-1$: $-1{,}5$. Ο $-\frac{1}{2}=-0{,}5$ είναι στη μέση του $-1$ και του $0$.\par
\apantisi $A(-3)$, $B(2{,}5)$, $\Gamma(-1{,}5)$, $\Delta(4)$ και ο $-0{,}5$ ανάμεσα στο $-1$ και στο $0$.\par
\elegxos Τα $-3$ και $-1{,}5$ είναι αριστερά του $0$, άρα αρνητικοί.

\tmima{Ασκήσεις}

\zoni{Το κατάλαβα}{πρόσημα, ακέραιοι, άξονας}

\ex Εξέφρασε με αριθμό: α) $5$ βαθμοί κάτω από το μηδέν \quad β) $12$ μ. πάνω από τη θάλασσα \quad γ) χρέος $8$ € \quad δ) κέρδος $15$ € \quad ε) το δεύτερο υπόγειο \quad στ) $30$ μ. κάτω από την επιφάνεια της θάλασσας.

\ex Ποιοι από τους αριθμούς $-7$, $+3{,}5$, $0$, $-\frac{1}{2}$, $8$, $+0{,}01$ είναι θετικοί, ποιοι αρνητικοί και ποιος δεν είναι ούτε το ένα ούτε το άλλο;

\ex Πες αν είναι ομόσημοι ή ετερόσημοι: α) $-3$ και $-8$ \quad β) $+4$ και $-4$ \quad γ) $0{,}5$ και $\frac{7}{2}$ \quad δ) $-2$ και $2{,}5$.

\ex Ποιοι από τους $-4$, $2{,}5$, $0$, $\frac{9}{3}$, $-\frac{1}{2}$, $10$ είναι ακέραιοι;

\ex Σχεδίασε άξονα με μονάδα $1$ cm και σημείωσε τους $-4$, $-1{,}5$, $0$, $2$, $3{,}5$.

\ex Στον άξονα βρες τις τετμημένες των $A,\,B,\,\Gamma,\,\Delta$.
\begin{center}
\akgr{-5}{5}{-4/A,-1.5/B,1/\Gamma,3.5/\Delta}
\end{center}

\ex Σωστό ή λάθος; α) Το $0$ είναι θετικός αριθμός. \quad β) Κάθε φυσικός αριθμός είναι ακέραιος. \quad γ) Κάθε ακέραιος είναι ρητός. \quad δ) Οι $-3$ και $-5$ είναι ετερόσημοι. \quad ε) Το $\frac{2}{3}$ είναι ρητός. \quad στ) Στον άξονα οι αρνητικοί βρίσκονται αριστερά του $0$.

\zoni{Το χρησιμοποιώ}{θερμόμετρα, ασανσέρ, θάλασσες}

\ex \textbf{Ο καιρός.} Οι προβλέψεις: Αθήνα $7$° πάνω από το $0$, Θεσσαλονίκη $3$° κάτω, Ιωάννινα $5$° κάτω, Χανιά $11$° πάνω, Φλώρινα $10$° κάτω. Γράψε τις θερμοκρασίες ως ρητούς και τοποθέτησέ τες σε έναν κατακόρυφο άξονα.

\ex \textbf{Ο ανελκυστήρας.} Οι όροφοι του ξενοδοχείου είναι $-2,\,-1,\,0,\,1,\,2,\,3$. Ο Νίκος είναι στον $3$ και θέλει το δεύτερο υπόγειο. Πόσους ορόφους κατεβαίνει;

\ex \textbf{Βουνό και θάλασσα.} Το αεροπλάνο πετάει στα $+300$ μ., το ψάρι στα $-45$ μ. και το υποβρύχιο στα $-120$ μ. Ποιο από τα τρία είναι πιο κοντά στην επιφάνεια της θάλασσας;

\ex \textbf{Το ταμείο.} Η Μυρτώ έβαλε $20$ €, ο Στέλιος πήρε $7$ €, η Φανή έβαλε $5$ € και ο Βασίλης πήρε $12$ €. Γράψε κάθε κίνηση ως ρητό.

\ex \textbf{Ο καιρός αλλάζει.} Το πρωί στο Μέτσοβο είναι $-3$° και μέχρι το μεσημέρι ανεβαίνει $5$ βαθμούς. Βρες τη νέα θερμοκρασία μετρώντας πάνω στον άξονα.

\zoni{Το σκέφτομαι}{είναι το μηδέν τίποτα;}

\ex Ο Στέλιος ισχυρίζεται ότι δεν υπάρχει αριθμός μικρότερος από το $0$. Φέρε αντιπαράδειγμα από τη ζωή.

\ex Γιατί το $0$ δεν είναι ούτε θετικός ούτε αρνητικός αριθμός;

\ex Τα $0$°C είναι η θερμοκρασία πήξης του νερού και τα $32$°F είναι η ίδια θερμοκρασία σε βαθμούς Φαρενάιτ. Τι δείχνει αυτό για το «μηδέν» στο θερμόμετρο;

\ex Πόσοι ακέραιοι υπάρχουν από το $-5$ ως το $5$, με τα άκρα; Πόσοι ανάμεσα στο $-5$ και στο $5$ χωρίς τα άκρα;

\zoni{Για το σπίτι}{ψάξε τους αρνητικούς γύρω σου}

\ex Βρες πέντε περιπτώσεις της καθημερινής ζωής που χρησιμοποιούνται αρνητικοί αριθμοί.

\ex Μέτρησε τη θερμοκρασία κάθε πρωί για μια εβδομάδα και παράστησέ τη σε άξονα.
